
\begin{ledgroupsized}[r]{120mm}
\footnotesize
\pstart
\noindent\textbf{Überlieferung:}
\pend
\end{ledgroupsized}
\begin{ledgroupsized}[r]{114mm}
\footnotesize
\pstart \parindent -6mm
\makebox[6mm][l]{\textit{L}}%
Aufzeichnung:
LH~XXXV~10, 14~Bl.~7. 
Ein Blatt~4\textsuperscript{o};
\href{\Scan{LH_35_10_14_007wz}}%
{Fragment einer Wasserzeichen-Gegenmarke} (Papier aus Osterode im Harz) 
am Blattrand;
oberer und linker Rand beschnitten. 
Eindrittel Seite auf Bl.~7~r\textsuperscript{o}; Bl.~7~v\textsuperscript{o} leer.
Der Träger bildete ursprünglich eine Einheit mit LH~XXXV~10, 14~Bl.~6 (Träger von N.\;\ref{60113}).
Am oberen Seitenrand, in Rot, der Vermerk von unklarer Hand \textit{VII}
(vgl. die Träger der benachbarten Signaturen LH~XXXV~10, 14 Bl.~2 bis 13, u.\,a. 
N.\;\ref{60117},
\ref{60118},
\ref{60115},
\ref{60112},
\ref{60116},
\ref{60113},
\ref{60110},
\ref{41824}
und \ref{60111}
in diesem Band).
\pend
\end{ledgroupsized}

\vspace{5mm}
\begin{ledgroup}
\footnotesize
\pstart
\noindent
\textbf{Datierungsgründe:}
Der Träger des vorliegenden Stücks N.\;\ref{60114} hing ursprünglich mit dem Quartblatt LH~XXXV~10, 14~Bl.~6 zusammen, auf dem das Konzept "Sphaera gravis sustentata a duobus planis" (N.\;\ref{60113}) überliefert ist. 
Zusammen bilden die Blätter etwa vier Fünftel eines Bogens, dessen verbleibenden Teile bislang nicht identifiziert werden konnten.
Der Träger von N.\;\ref{60113} enthält ein Wasserzeichen mit dem Motiv "Wappen der Stadt Osterode\indexO{Osterode} (Harz)"; N.\;\ref{60114} enthält die Gegenmarke; das Zeichen ist für den Zeitraum 1690 bis 1694 belegt, also für die Zeit unmittelbar nach der italienischen Reise (wobei Leibniz erst in der zweiten Junihälfte 1690 in Hannover eintraf).
Damit ist der Terminus post quem für die Entstehung von N.\;\ref{60114} gesichert.
\pend
\pstart
Der Vergleich mit dem Konzept "Solutio problematis R.\,P.\;Vannii S.\,J." von April bis Juni 1690 (siehe dazu die Einleitung zu N.\;\ref{intro_vanni-solutio-problematis})
und der Aufzeichnung "Arcanum eccentrici concursus detectum" (10. November 1690 -- September 1691, N.\;\ref{60339})
liefert einen wahrscheinlichen Terminus ante quem.
Die Aufzeichung N.\;\ref{60114} führt denselben Gedanken auf, außerdem mit wörtlichen Anklängen, wie die Passage auf
\refpassage{35_10_13_010-011_14a}{35_10_13_010-011_14b}
des umfangreicheren beider Zeugen der "Solutio" (\textit{L\textsuperscript{2}} von N.\;\ref{41823}).
Bei N.\;\ref{60114} handelt es sich vermutlich um einen Versuch, das dort festgestellte Ergebnis zu präzisieren oder separat festzuhalten.
Derselbe Gedanke kommt wiederum in der Aufzeichnung "Arcanum" zur Anwendung (\refpassage{LH_37_05_124-125_125r_conatuseccentr-1}{LH_37_05_124-125-perinde}).
Aus diesen Gründen kann von einer wahrscheinlichen Entstehung von N.\;\ref{60114} noch in zeitlicher Nähe zu N.\;\ref{41823}, womöglich in Überschneidung mit der Abfassung von  N.\;\ref{60339}, ausgegangen werden, woraus sich die vorgeschlagene Datierung ergibt.
\pend 
\end{ledgroup}
\clearpage
\selectlanguage{latin}
\vspace{8mm}
\pstart
\noindent
\href{\Scan{LH_35_10_14_007r}}%
{\lbrack7~r\textsuperscript{o}\rbrack}\ 
\pend 
\begin{edfigure}[0.25em]%
\hspace{-2em}%
\includegraphics[scale=1.0]{%
gesamttex/edit_VIII-4/images/LH_35_10_14_007_d.pdf}
\edcaption[0.25em]{%
\lbrack\textit{Fig.~1}\rbrack}
\end{edfigure}
\pstart
Grave\indexS{grave} \textit{GAB},
\edtext{cujus centrum grav.\indexS{centrum gravitatis} sit \textit{G},}{%
\lemma{}%
\Bfootnote{%
cujus centrum grav. sit \textit{G}, 
\textit{erg.~L}}}
incumbit\indexS{grave duobus planis incumbens}
\edtext{planis \textit{AC}}{%
\lemma{planis}%
\Bfootnote{%
\textit{(1)}~\textit{A}
\textit{(2)}~\textit{AC}~\textit{L}}}
et \textit{BD}. Jungantur rectae \textit{AG}, \textit{BG}. Perinde est ac si grave\indexS{grave} punctum \textit{G}
per contos\indexS{contus} \textit{GA}, \textit{GB} innitatur 
\edtext{planis\indexS{planum} vel}{%
\lemma{planis}%
\Bfootnote{%
\textit{(1)}~. Sed conti impellantur
\textit{(a)}~norma
\textit{(b)}~in normalibus ad plana
\textit{(2)}~vel~\textit{L}}}
vicissim, mobile \textit{G}, sublata consideratione gravitatis,\indexS{gravitas} impellatur 
per contos\indexS{contus} \textit{AG}, \textit{BG}\lbrack;\rbrack\
intelligi autem debet impulsum\indexS{impulsus} fieri in perpendicularibus ad \textit{AC}\bracks{,} 
\edtext{\bracks{\textit{BD}}.}{\lemma{}\Bfootnote{\textit{BD} \textit{erg.~Hrsg.}}}
Perinde ac 
\edtext{si recta}{%
\lemma{si}%
\Bfootnote{%
\textit{(1)}~\textit{AE}
\textit{(2)}~planum
\textit{(a)}~recta~\textit{L}}}
\edtext{\textit{A}\lbrack\textit{CE}\rbrack}{%
\lemma{}%
\Bfootnote{%
\textit{AEC} 
\textit{L ändert Hrsg.}}}
moveatur circa \textit{E} et recta \textit{FBD} circa \textit{F}\lbrack,\rbrack\
quo pacto constat perinde esse ac 
\edtext{si mobile}{%
\lemma{si}%
\Bfootnote{%
\textit{(1)}~grave
\textit{(2)}~mobile~\textit{L}}}
impellatur nisu \textit{CG} vel \textit{DG}. %
\pend
\pstart
Itaque perinde est ac si \textit{G} impellatur in rectis\indexS{recta normalis ad planum} normalibus ad plana.\indexS{planum}
\pend
\pstart
\edtext{Et vis\indexS{vis}}{%
\lemma{Et}%
\Bfootnote{%
\textit{(1)}~pr
\textit{(2)}~ad aequili
\textit{(3)}~vis~\textit{L}}}
\textit{GS} premet plano\indexS{vis planum premens} gravitationibus\indexS{gravitatio} \textit{GR}, \textit{GQ}.
\pend
